\begin{textsample}{NEG Dim 1 – 2022\_c4 – Score -4.00 – 2022\_c4/enade\_2022\_tec\_gastr\_1935a4d527.txt}  \label{ex:f1_neg_174}
CAPÍTULO 5 PERCEPÇÃO DA PROVA Nas análises feitas neste capítulo, constam as percepções dos concluintes da Área de Tecnologia em Gastronomia sobre a prova aplicada no Enade/2022.
Essas percepções foram mensuradas por meio de nove questões que avaliaram desde o grau de dificuldade da prova até o tempo gasto para concluí -la.
As percepções sobre a prova foram relacionadas ao desempenho dos estudantes, à grande região de funcionamento do curso, à categoria administrativa e ao tipo de organização acadêmica da IES.
O questionário de percepção da prova encontra -se ao final do Anexo VIII, no qual \textbf{está} a reprodução do exame.
Na apresentação dos dados relativos às nove questões sobre as percepções a respeito da prova, o desempenho dos estudantes foi classificado em quatro quartos.
Para tanto, esse desempenho foi ordenado de forma ascendente.
O percentil 25, P25, também conhecido como primeiro quartil, é a nota de desempenho que deixa um quarto ( 25 \% ) dos valores observados abaixo e três quartos acima.
O quarto inferior de desempenho é composto pelas notas abaixo do primeiro quartil.
Já o percentil 75, P75, também conhecido como terceiro quartil, é o valor para o qual há três quartos ( 75 \% ) dos dados abaixo e um quarto acima dele.
O quarto superior de desempenho é composto pelas notas iguais ou acima do terceiro quartil.
O percentil 50, P50, também conhecido como mediana, é o valor que divide as notas em dois conjuntos de igual tamanho.
O segundo quarto inclui os valores entre o primeiro quartil ( P25 ) e a mediana.
O terceiro quarto contém os valores entre a mediana ( P50 ) e o terceiro quartil ( P75 ).
Vale ressaltar que percentis, quartis e medianas são pontos que, não obrigatoriamente, pertencem ao conjunto original de dados, ao passo que os quartos são subconjuntos dos dados originais.
Na Figura 5.1, apresenta -se uma ilustração dos quatro quartos descritos.
A seguir, serão apresentados gráficos com resultados selecionados, relativos às nove questões avaliadas por grupos de estudantes.
Nas barras dos gráficos, apresenta -se o percentual de estudantes que assinalaram uma das opções ou a soma das porcentagens daqueles que assinalaram duas ( ou três ) delas.
Por exemplo, para as questões 1 e 2, nos gráficos, é apresentada a porcentagem total de participantes que assinalaram as opções ( D ) “ Difícil ” e ( E ) “ Muito difícil ”.
Em \textbf{cada} barra, foram assinalados, também, os extremos do intervalo de confiança de 95 \% como linhas verticais unidas por uma linha horizontal na forma da letra H maiúscula.
O estimador de um parâmetro com um \textbf{certo} nível de confiança ( por exemplo, 95,0 \% ) deve conter o parâmetro no intervalo de confiança em 95 \% das vezes.
Na comparação entre os estimadores dos parâmetros de duas classes de uma dada categoria ( por exemplo : Norte e Nordeste nas Grandes Regiões, de primeiro e último quarto dentro de desempenho, de IES públicas e privadas, ou de universidades e faculdades ), associados aos seus respectivos intervalos de confiança, diz -se que não há diferença estatisticamente significativa entre os parâmetros das duas categorias quando há interseção entre os intervalos de confiança, e que há diferença, se os intervalos de confiança são disjuntos ( para maiores informações vide Glossário ).
Nas Tabelas do Anexo V, são apresentados os valores absolutos e a distribuição percentual10 das alternativas válidas das nove questões, segundo o mesmo recorte de desempenho dos estudantes e da grande região de funcionamento do curso, de categoria administrativa e do tipo de organização acadêmica da IES. 5.1 GRAU DE DIFICULDADE DA PROVA Nesta seção são apresentadas as percepções dos estudantes quanto ao grau de dificuldade da prova Enade/2022, representadas nos gráficos seguintes, que contemplam o Componente de Formação Geral e o Componente Específico. 10 Reitera -se que uma das convenções para tabelas numéricas ( Subseção 2.2.3 do Capítulo 2 ) refere -se à possibilidade de, por questão de arredondamento, a soma das partes não resultar em 100 \%. 5.1.1 COMPONENTE DE FORMAÇÃO GERAL O Gráfico 5.1 apresenta as respostas obtidas para a Questão 1 : “ Qual o grau de dificuldade desta prova na parte de Formação Geral? ”.
De acordo com os dados do Gráfico 5.1, a região Norte apresentou o maior percentual de estudantes que consideraram a prova “ Difícil ” ou “ Muito Difícil ”, com 29,6 \%, seguida da região Sul com 23 \%, a região Centro-Oeste com 22,8 \%, a região Sudeste com 19,4 \% e a região Nordeste 15,5 \%.
O Gráfico 5.2 traz mais resultados sobre a questão tratada no gráfico anterior, porém por quartos de desempenho.
O percentual de estudantes que consideraram o componente de Formação Geral da prova “ Difícil ” ou “ Muito difícil ” foi decrescente em função dos quartos de desempenho : 30,1 \%, no primeiro quarto, e 12,6 \%, no quarto, o grupo de melhor desempenho na prova.
Nos quartos de desempenho intermediários, a proporção de alunos que consideraram a prova “ Difícil ” ou “ Muito difícil ” foi de 24 \%, no segundo quarto, e de 20,5 \%, no terceiro quarto.
O Gráfico 5.3 traz outros resultados também em relação à mesma questão retratada no gráfico anterior, porém por organização acadêmica das IES.
O percentual de estudantes que consideraram o componente de Formação Geral da prova “ Difícil ” ou “ Muito difícil ” foi maior entre os estudantes de universidades, alcançando 24,8 \%, enquanto entre os estudantes dos CEFET/IF, a organização acadêmica com a menor incidência, o percentual alcançou 15,7 \%.
Nos centros universitários, o percentual foi de 18,8 \% e nas faculdades, 18,6 \%.
O Gráfico 5.4 apresenta os dados por categoria administrativa das IES.
O percentual de estudantes que consideraram o componente de Formação Geral da prova “ Difícil ” ou “ Muito difícil ” foi maior entre os estudantes de cursos de IES privadas ( 22,1 \% ). 5.1.2 COMPONENTE DE CONHECIMENTO ESPECÍFICO O Gráfico 5.5 apresenta as respostas obtidas para a Questão 2 : “ Qual o grau de dificuldade desta prova na parte de Componente Específico? ”.
A análise das respostas dos estudantes quanto ao grau de dificuldade do Componente Específico da prova, agregado por grande região, mostra similaridades estatísticas, em que o percentual mais alto foi encontrado na região Norte ( 22,7 \% ) e mais baixo na região Sul ( 17,2 \% ).
As demais regiões apresentaram os seguintes percentuais : Centro-Oeste ( 22 \% ), Nordeste ( 20 \% ) e Sudeste ( 17,8 \% ).
O Gráfico 5.6 traz mais resultados sobre a questão tratada no gráfico anterior, porém por quartos de desempenho.
Considerando -se a avaliação da dificuldade das questões do Componente Específico da prova de acordo com o desempenho dos estudantes, observa -se uma tendência não linear dos que consideraram a parte específica “ Difícil ” ou “ Muito difícil ”, com o terceiro quarto com 20,3 \% e o segundo quarto com 17,1 \%, contrariando a ideia de que a taxa de declaração “ Difícil ” ou “ Muito Difícil ” deveria ser maior no segundo quarto.
Nos dois extremos temos, 26,4 \% no primeiro quarto e 14,3 \% no quarto.
O Gráfico 5.7 traz outros resultados também em relação à mesma questão retratada no gráfico anterior, porém por organização acadêmica das IES.
A análise das respostas dos estudantes quanto ao grau de dificuldade do Componente Específico da prova, agregado por organização acadêmica das IES, mostra que o percentual mais alto \textbf{está} nos CEFET/IF ( 24 \% ) e o percentual mais baixo \textbf{está} nas faculdades ( 16,3 \% ).
As demais organizações acadêmicas apresentam os seguintes percentuais : centros universitários com 17,1 \% e universidades com 20,7 \%.
O Gráfico 5.8 apresenta os dados por categoria administrativa das IES.
A análise das respostas dos estudantes quanto ao grau de dificuldade do Componente Específico da prova, agregado por categoria administrativa, mostra que as instituições públicas apresentaram um percentil maior ( 22,9 \% ), quando comparadas com as instituições privadas ( 18,9 \% ). 5.2 EXTENSÃO DA PROVA EM RELAÇÃO AO TEMPO TOTAL O Gráfico 5.9 apresenta as respostas obtidas para a Questão 3 : “ Considerando a extensão da prova, em relação ao tempo total, você considera que a prova foi ”.
A análise das respostas dos estudantes quanto ao grau de extensão da prova como “ Longa ” ou “ Muito longa ”, agregado por grande região, mostra similaridades estatísticas, em que o percentual mais alto foi encontrado na região Norte ( 35,6 \% ) e mais baixo na região Nordeste ( 31,7 \% ).
O Gráfico 5.10 traz mais resultados sobre a questão tratada no gráfico anterior, porém por quartos de desempenho.
A análise das respostas dos estudantes quanto ao grau de extensão da prova como “ Longa ” ou “ Muito longa ”, agregado por quartos de desempenho, mostra uma tendência linear entre a extensão da prova e o quarto de desempenho.
O 1{\EmojiFont °} Quarto apresenta o maior percentil de estudantes que classificaram a prova como “ Longa ” ou “ Muito longa ”, com 40,8 \% ; na outra extremidade, o 4{\EmojiFont °} Quarto apresentou o \textbf{menor} percentil, com 27,5 \%.
O Gráfico 5.11 traz outros resultados também em relação à mesma questão retratada no gráfico anterior, porém por organização acadêmica das IES.
A análise das respostas dos estudantes quanto à extensão da prova em relação ao tempo total, agregado por organização acadêmica das IES, mostra similaridade de percentuais, com 35,6 \% nas faculdades, 35,3 \% nas universidades, 33,1 \% nos centros universitários e 31,5 \% nos CEFET/IF.
O Gráfico 5.12 apresenta os dados por categoria administrativa das IES.
A análise das respostas dos estudantes quanto à extensão da prova, agregada por categoria administrativa, mostra que as instituições privadas apresentaram um percentil levemente maior ( 34,5 \% ), quando comparada com as instituições públicas ( 32,4 \% ). 5.3 COMPREENSÃO DOS ENUNCIADOS DAS QUESTÕES Também organizados em relação ao Componente de Formação Geral e ao Componente Específico, os dados apresentados nesta seção dizem respeito à forma como os estudantes compreenderam os enunciados das questões da prova, considerando as formulações expressas nesses mesmos enunciados. 5.3.1 COMPONENTE DE FORMAÇÃO GERAL O Gráfico 5.13 apresenta as respostas obtidas para a Questão 4 : “ Os enunciados das questões da prova na parte de Formação Geral \textbf{estavam} claros e objetivos? ”.
A análise das respostas dos estudantes quanto à clareza do enunciado das questões da prova de Formação Geral, agregado por grande região, mostra similaridades estatísticas, em que o percentual mais alto foi encontrado na região Centro-Oeste ( 83,3 \% ) e mais baixo na região Norte ( 76,3 \% ).
O Gráfico 5.14 traz mais resultados sobre a questão tratada no gráfico anterior, porém por quartos de desempenho.
A análise das respostas dos estudantes quanto ao grau de clareza dos enunciados das questões da prova de Formação Geral, agregado por quartos de desempenho, mostra uma tendência linear entre o entendimento dos enunciados como “ claros e objetivos ” e o quarto de desempenho.
O 1{\EmojiFont °} Quarto apresenta o \textbf{menor} percentil de estudantes que classificaram a prova como contendo enunciados “ claros e objetivos ”, com 68,3 \% ; na outra extremidade, o 4{\EmojiFont °} Quarto apresentou o maior percentil, com 89 \%.
O Gráfico 5.15 traz outros resultados também em relação à mesma questão retratada no gráfico anterior, porém por organização acadêmica das IES.
A análise das respostas dos estudantes quanto ao grau de clareza dos enunciados das questões da prova de Formação Geral, agregado por organização acadêmica das IES, mostra percentuais semelhantes, com 82,5 \% nos CEFET/IF, 77,8 \% nas universidades, 79,4 \% nos centros universitários e 79,3 \% nas faculdades.
O Gráfico 5.16 apresenta os dados por categoria administrativa das IES.
A análise das respostas dos estudantes quanto ao grau de clareza dos enunciados das questões da prova de Formação Geral, agregado por categoria administrativa, mostra que as instituições públicas apresentaram um percentil maior ( 82,8 \% ) de entendimento dos enunciados, quando comparada com as instituições privadas ( 78,4 \% ). 5.3.2 COMPONENTE ESPECÍFICO O Gráfico 5.17 apresenta as respostas obtidas para a Questão 5 : “ Os enunciados das questões da prova na parte de Componente Específico \textbf{estavam} claros e objetivos? ”.
A análise das respostas dos estudantes quanto à clareza do enunciado das questões da prova de Conhecimento Específico, agregado por grande região, demonstra que as cinco regiões apresentaram percentuais próximos, variando de 81,4 \% na região Nordeste até 88,3 \% na região Sul.
O Gráfico 5.18 traz mais resultados sobre a questão tratada no gráfico anterior, porém por quartos de desempenho.
A análise das respostas dos estudantes quanto ao grau de clareza dos enunciados das questões da prova de Conhecimento Específico, agregado por quartos de desempenho, mostra uma tendência linear entre o entendimento dos enunciados como claros e objetivos e o quarto de desempenho.
O 1{\EmojiFont °} Quarto apresenta o \textbf{menor} percentil de estudantes que classificaram a prova como contendo enunciados claros e objetivos, com 71,7 \% e na outra extremidade, o 4{\EmojiFont °} Quarto, apresentou o maior percentil, com 92,4 \%.
O Gráfico 5.19 traz outros resultados também em relação à mesma questão retratada no gráfico anterior, porém por organização acadêmica das IES.
A análise das respostas dos estudantes quanto ao grau de clareza dos enunciados das questões da prova de Conhecimento Específico, agregado por organização acadêmica das IES, mostra percentuais semelhantes, com 88,3 \% nas faculdades, 84,8 \% nos centros universitários, 83,2 \% nas universidades e 83 \% nos CEFET/IF.
O Gráfico 5.20 apresenta os dados por categoria administrativa das IES.
A análise das respostas dos estudantes quanto ao grau de clareza dos enunciados das questões da prova de Conhecimento Específico, agregado por categoria administrativa, mostra que as instituições privadas apresentaram um percentil levemente maior ( 84,3 \% ) de entendimento dos enunciados, quando comparada com as instituições públicas ( 83,4 \% ). 5.4 SUFICIÊNCIA DAS INFORMAÇÕES/INSTRUÇÕES FORNECIDAS O Gráfico 5.21 apresenta as respostas obtidas para a Questão 6 : “ As informações/instruções fornecidas para a resolução das questões foram suficientes para resolvê -las? ”.
As análises das respostas dos estudantes quanto ao grau de instruções fornecidas foram suficientes em todas ou na maioria das questões, agregado por grande região, demonstra que as cinco regiões apresentaram percentuais próximos com 83,7 \% nas regiões Norte e Nordeste, 86,4 \% na região Sudeste, 87,4 \% na região Centro-Oeste e 89,3 \% na região Sul.
O Gráfico 5.22 traz resultados sobre a questão tratada no gráfico anterior, porém por quartos de desempenho.
A análise das respostas dos estudantes quanto ao grau instruções fornecidas foram suficientes em todas ou na maioria das questões, agregado por quartos de desempenho, mostra uma tendência linear entre as instruções fornecidas sendo suficientes e o quarto de desempenho.
O 1{\EmojiFont °} Quarto apresenta o \textbf{menor} percentil, com 76,7 \% e, na outra extremidade, o 4{\EmojiFont °} Quarto apresentou o maior percentil, com 93,6 \%.
O Gráfico 5.23 traz outros resultados também em relação à mesma questão retratada no gráfico anterior, porém por organização acadêmica das IES.
Na análise por organização acadêmica, a proporção de respondentes que escolheram a opção em todas ou na maioria das questões “ informações/instruções fornecidas para a resolução das mesmas ” foi de 87,7 \% nos CEFET/IF, 86,6 \% nos centros universitários e faculdades, e 85,3 \% nas universidades.
O Gráfico 5.24 apresenta os dados por categoria administrativa das IES.
Em relação à categoria administrativa, observa -se que, nas IES Públicas, 87,4 \% ofereceram como resposta que em todas ou na maioria das questões as instruções fornecidas para a resolução das mesmas eram suficientes e, nas IES privadas, o percentual foi de 85,9 \%. 5.5 DIFICULDADE ENCONTRADA AO RESPONDER À PROVA O Gráfico 5.25 apresenta as respostas obtidas para a Questão 7 : “ Você se deparou com alguma dificuldade ao responder à prova?
Qual? ”.
Nos Gráficos 5.25 a 5.28, apresentados de forma sequenciada mais adiante, são expostos os percentuais de estudantes que apontaram o desconhecimento do conteúdo como “ dificuldade ao responder à prova ”.
O Gráfico 5.25 apresenta as respostas obtidas sobre o desconhecimento do conteúdo como a principal “ dificuldade ao responder à prova ”.
Considerando -se as cinco grandes regiões brasileiras, os percentuais dos que utilizaram desconhecimento do conteúdo como a principal “ dificuldade ao responder à prova ” foram : Sul ( 23,1 \% ), Nordeste ( 20,1 \% ), Sudeste ( 20,1 \% ), Norte ( 19,7 \% ) e Centro-Oeste ( 17,1 \% ).
O Gráfico 5.26 traz resultados sobre o ponto tratado no gráfico anterior, porém por quartos de desempenho.
A análise das respostas dos estudantes quanto ao desconhecimento do conteúdo como a principal dificuldade ao responder à prova, agregado por quartos de desempenho, mostra valores não lineares : no 1{\EmojiFont °} Quarto tem -se 22,4 \% ; 2{\EmojiFont °} Quarto, 20,2 \% ; 3{\EmojiFont °} Quarto, 19,3 \% ; e 4{\EmojiFont °} Quarto, 19,7 \%.
O Gráfico 5.27 traz outros resultados também em relação à mesma questão retratada no gráfico anterior, porém por organização acadêmica das IES.
Na análise por organização acadêmica, a proporção de respondentes que escolheram a opção desconhecimento do conteúdo como a principal dificuldade ao responder à prova variou de 25,6 \% ( CEFET/IF ) a 19,4 \% ( universidades ).
O Gráfico 5.28 apresenta os dados por categoria administrativa das IES.
Em relação à categoria administrativa, observa -se que, nas IES públicas, 23,2 \% ofereceram como resposta desconhecimento do conteúdo como a principal dificuldade ao responder à prova e nas IES privadas o percentual foi de 20 \%. 5.6 CONTEÚDOS DAS QUESTÕES OBJETIVAS DA PROVA O Gráfico 5.29 apresenta as respostas obtidas para a Questão 8 : “ Considerando apenas as questões objetivas da prova, você percebeu que não estudou ainda a maioria desses conteúdos? ”.
Considerando -se as cinco grandes regiões brasileiras, os percentuais dos que utilizaram não estudaram ainda a maioria desses conteúdos foram : Centro-Oeste ( 7,8 \% ), Nordeste ( 7,8 \% ), Norte ( 6 \% ), Sul ( 5 \% ) e Sudeste ( 4,7 \% ).
O Gráfico 5.30 traz resultados sobre a questão tratada no gráfico anterior, porém por quartos de desempenho.
Uma vez analisadas as alternativas, escolhidas pelos estudantes que se situam nos diferentes quartos de desempenho, observa -se uma tendência crescente : uma proporção menor de participantes no quarto superior ( 2,1 \% ) declarou que “ não estudaram ainda a maioria desses conteúdos ” quando comparada com a proporção de participantes nos quartos inferiores, especialmente os do 1º Quarto ( 11,4 \% ).
O Gráfico 5.31 traz outros resultados também em relação à mesma questão retratada no gráfico anterior, porém por organização acadêmica das IES.
Na análise por organização acadêmica, a proporção de respondentes que escolheu a opção “ não estudaram ainda a maioria desses conteúdos ” foi pequena.
Observa -se 7,1 \% nos CEFET/IF ; 5,8 \% nas universidades ; 5,7 \% nos centros universitários e 5,3 \% nas faculdades.
O Gráfico 5.32 apresenta os dados por categoria administrativa das IES.
Em relação à categoria administrativa, observa -se que, nas IES públicas, 8 \% ofereceram como resposta que “ não estudaram ainda a maioria desses conteúdos ” e, nas IES privadas, o percentual foi de 5,6 \%. 5.7 TEMPO GASTO PARA CONCLUIR A PROVA O Gráfico 5.33 apresenta as respostas obtidas para a Questão 9 : “ Qual foi o tempo gasto por você para concluir a prova? ”.
Considerando -se as cinco grandes regiões brasileiras, os percentuais dos que utilizaram entre duas e quatro horas para finalizar a prova foram : Nordeste ( 86,5 \% ), Norte ( 76,4 \% ), Sul ( 75,2 \% ), Centro-Oeste ( 72,3 \% ) e Sudeste ( 71,5 \% ).
O Gráfico 5.34 traz resultados sobre a questão tratada no gráfico anterior, porém por quartos de desempenho.
Uma vez analisadas as alternativas, escolhidas pelos estudantes que se situam nos diferentes quartos de desempenho, observa -se uma tendência crescente : uma proporção maior de participantes no quarto superior ( 81,8 \% ) declarou ter gastado entre duas e quatro horas para concluir a prova quando comparada com a proporção de participantes nos quartos inferiores, sendo o primeiro quarto com 69,6 \%.
O Gráfico 5.35 traz outros resultados também em relação à mesma questão retratada no gráfico anterior, porém por organização acadêmica das IES.
Na análise por organização acadêmica, os percentuais dos que utilizaram entre duas e quatro horas para finalizar a prova teve a seguinte distribuição : CEFET/IF ( 86,6 \% ), faculdades ( 82,1 \% ), centros universitários ( 75,3 \% ) e universidades ( 73,5 \% ).
O Gráfico 5.36 apresenta os dados por categoria administrativa das IES.
O Gráfico 5.36 permite inferir que o percentual de estudantes das IES públicas ( 83,5 \% ) que perceberam que gastar de duas a quatro horas para concluir a prova foi superior ao de estudantes das IES privadas ( 75,2 \% ).
Os resultados do Enade 2022, apresentados neste capítulo, contribuem para o entendimento de que a consideração da percepção dos estudantes concluintes sobre a prova do Enade é importante tanto para a compreensão do seu desempenho como para a reflexão sistemática e constante a respeito de determinadas realidades relacionadas a esse mesmo desempenho, no país.
Por outro lado, como todos os elementos do trabalho pedagógico institucional mostram -se diretamente relacionados com o desempenho dos estudantes, também ganha relevância conhecer e analisar os resultados da percepção dos coordenadores de curso sobre a mesma prova.
As informações expressas por estudantes e coordenadores sobre questões diversas associadas aos dois componentes da prova – Formação Geral e Conhecimento Específico – podem agregar valor em ações institucionais voltadas para o aperfeiçoamento do instrumento avaliativo, considerando os vários elementos envolvidos em sua elaboração.
Neste contexto, é relevante que os resultados das percepções de estudantes e coordenadores sobre a prova do Enade ganhem peso como um elemento a mais nas análises da oferta dos cursos nas modalidades presencial e a distância, em nível nacional.
Espera -se, assim, que, em uma visão formativa de avaliação, o conhecimento e a análise desses resultados possam vir a contribuir para a construção de alternativas institucionais que venham a induzir a qualidade da educação superior.

% matched lemmas: cada, certo, estar, menor
\end{textsample}
